% ==========================================================
% 030_slice_domains.tex
% Forecast Admissibility Surface (FAS)
% ==========================================================

\section{Slice Domains and Admissibility Granularity}
\label{sec:fas_slice_domains}

The Forecast Admissibility Surface operates over a \emph{slice domain} that
partitions the full forecast space into structurally distinct units of analysis.
A slice is defined by a fixed set of keys—such as entity, time-of-day interval, or
site—that jointly determine the granularity at which admissibility decisions are
made. The choice of slice domain is consequential: admissibility is not an
intrinsic property of a demand process in the abstract, but a property of that
process \emph{at a particular resolution and conditioning}.

Let $\Domain$ denote the full forecast domain under consideration, comprising all
entities, sites, and time indices eligible for forecasting. FAS induces a
partition of $\Domain$ into disjoint slices
$\{\Slice_k\}_{k \in \SliceKeys}$, where each slice $\Slice_k$ corresponds to a
unique realization of a slice key tuple. All admissibility decisions are made at
the level of these slices.

\subsection{Supported Slice Keys}

FAS is agnostic to the specific choice of slice keys, provided that the resulting
slices admit sufficient observational support for baseline error analysis. In
practice, three slice granularities are commonly employed:

\begin{itemize}
  \item \textbf{Entity-level slices}, in which all observations for a given
        forecast entity are aggregated across time and location.
  \item \textbf{Entity-by-interval slices}, which condition admissibility on both
        the forecast entity and the evaluation interval (e.g., half-hour of day),
        allowing time-of-day–specific hostility to be detected.
  \item \textbf{Site-by-entity-by-interval slices}, which further condition on
        location, capturing localized executional or demand-driven pathologies.
\end{itemize}

These slice modes correspond to increasingly fine-grained admissibility surfaces.
Coarser slices provide greater statistical support but may obscure localized
instability, while finer slices offer greater specificity at the cost of reduced
support.

\subsection{Granularity and Support Tradeoffs}

Admissibility classification is meaningful only insofar as each slice contains
sufficient observational mass. Extremely fine slice definitions can produce
spurious classifications driven by sampling noise rather than structural
behavior. Conversely, overly coarse slices may average over heterogeneous regimes
and falsely appear admissible.

FAS therefore treats slice granularity as a governance choice rather than a
learned parameter. Minimum support thresholds act as a conservative guardrail:
slices with insufficient observations are not declared admissible, but instead
classified as \Conditional to reflect uncertainty rather than hostility. This
ensures that lack of data is not conflated with structural inadmissibility.

\subsection{Joinability and Determinism}

A defining property of the slice domain in FAS is \emph{joinability}. Each slice
is defined using keys that exist natively in forecast panels, evaluation tables,
and governance artifacts. As a result, the admissibility surface produced by FAS
can be deterministically joined back onto downstream datasets at the same
granularity, constraining training, evaluation, and readiness control without
introducing ambiguity or additional state.

By grounding admissibility in explicit slice domains, FAS ensures that decisions
about where forecasting is permitted are transparent, reproducible, and
consistent across the Forecast Readiness Framework.
